\hypertarget{chapter-9-kegg-compound-and-reaction-integration}{%
\chapter{Chapter 9: KEGG Compound and Reaction
Integration}\label{chapter-9-kegg-compound-and-reaction-integration}}

\hypertarget{introduction}{%
\section{9.1 Introduction}\label{introduction}}

\textbf{Biochemical formula tracking} (Chapter 6) requires accurate
\textbf{elemental compositions} for all metabolites. Manually entering
formulas for hundreds of compounds is \textbf{error-prone and tedious}.
\textbf{This chapter presents automatic formula retrieval} from the
\textbf{Kyoto Encyclopedia of Genes and Genomes (KEGG)}, the largest
public database of biochemical compounds and reactions.

\textbf{KEGG database} (www.kegg.jp): - \textbf{KEGG COMPOUND}: 18,000+
metabolites with formulas, structures, pathways - \textbf{KEGG
REACTION}: 11,000+ biochemical reactions with stoichiometry, EC numbers
- \textbf{REST API}: HTTP endpoints for programmatic access (no
authentication required)

\textbf{SHYpn KEGG integration}: 1. \textbf{Compound lookup}: Fetch
formula by KEGG ID (e.g., ``C00031'' \textrightarrow ``C₆H₁₂O₆'') 2. \textbf{Reaction
parsing}: Extract stoichiometry from KEGG reaction equations 3.
\textbf{Pathway import}: Bulk import entire KEGG pathways as Bio-PN
networks 4. \textbf{Automatic enrichment}: Suggest missing cofactors
(ATP, NAD⁺, H₂O)

\textbf{Benefits}: - \textbf{Speed}: Entire glycolysis pathway (10
reactions) imported in \textless5 seconds - \textbf{Accuracy}: KEGG
formulas are manually curated (95\%+ accuracy) - \textbf{Completeness}:
Automatically adds cofactors often omitted in textbook representations

\begin{center}\rule{0.5\linewidth}{0.5pt}\end{center}

\hypertarget{kegg-rest-api}{%
\section{9.2 KEGG REST API}\label{kegg-rest-api}}

\hypertarget{endpoints}{%
\subsection{9.2.1 Endpoints}\label{endpoints}}

\textbf{Base URL}: \passthrough{\lstinline!https://rest.kegg.jp/!}

\textbf{Key endpoints} (GET requests):

\begin{longtable}[]{@{}
  >{\raggedright\arraybackslash}p{(\columnwidth - 4\tabcolsep) * \real{0.3125}}
  >{\raggedright\arraybackslash}p{(\columnwidth - 4\tabcolsep) * \real{0.4062}}
  >{\raggedright\arraybackslash}p{(\columnwidth - 4\tabcolsep) * \real{0.2812}}@{}}
\toprule\noalign{}
\begin{minipage}[b]{\linewidth}\raggedright
Endpoint
\end{minipage} & \begin{minipage}[b]{\linewidth}\raggedright
Description
\end{minipage} & \begin{minipage}[b]{\linewidth}\raggedright
Example
\end{minipage} \\
\midrule\noalign{}
\endhead
\bottomrule\noalign{}
\endlastfoot
\passthrough{\lstinline!get/<compound\_id>!} & Fetch compound details &
\passthrough{\lstinline!get/C00031!} \textrightarrow Glucose \\
\passthrough{\lstinline!get/<reaction\_id>!} & Fetch reaction details &
\passthrough{\lstinline!get/R00299!} \textrightarrow Hexokinase \\
\passthrough{\lstinline!find/compound/<query>!} & Search compounds &
\passthrough{\lstinline!find/compound/glucose!} \\
\passthrough{\lstinline!link/compound/<pathway\_id>!} & List compounds
in pathway & \passthrough{\lstinline!link/compound/map00010!} \textrightarrow
Glycolysis \\
\passthrough{\lstinline!list/pathway!} & List all pathways & Returns
500+ pathways \\
\end{longtable}

\textbf{Response format}: Plain text (tab-delimited fields)

\hypertarget{example-fetch-glucose}{%
\subsection{9.2.2 Example: Fetch Glucose}\label{example-fetch-glucose}}

\textbf{Request}:

\begin{lstlisting}
GET https://rest.kegg.jp/get/C00031
\end{lstlisting}

\textbf{Response}:

\begin{lstlisting}
ENTRY       C00031                      Compound
NAME        D-Glucose;
            Grape sugar;
            Dextrose
FORMULA     C6H12O6
EXACT_MASS  180.0634
MOL_WEIGHT  180.1559
REMARK      Same as: D00009
REACTION    R00299 R00300 R00301 R00302 R00303
PATHWAY     map00010  Glycolysis / Gluconeogenesis
            map00030  Pentose phosphate pathway
            map00052  Galactose metabolism
ENZYME      2.7.1.1   2.7.1.2   3.2.1.20
DBLINKS     PubChem: 3333
            ChEBI: 17234
            KNApSAcK: C00001161
///
\end{lstlisting}

\textbf{Key fields}: - \passthrough{\lstinline!FORMULA!}: C₆H₁₂O₆ (Hill
notation) - \passthrough{\lstinline!REACTION!}: List of reactions
involving this compound - \passthrough{\lstinline!PATHWAY!}: Associated
pathways - \passthrough{\lstinline!DBLINKS!}: Cross-references to
PubChem, ChEBI

\hypertarget{example-fetch-hexokinase-reaction}{%
\subsection{9.2.3 Example: Fetch Hexokinase
Reaction}\label{example-fetch-hexokinase-reaction}}

\textbf{Request}:

\begin{lstlisting}
GET https://rest.kegg.jp/get/R00299
\end{lstlisting}

\textbf{Response}:

\begin{lstlisting}
ENTRY       R00299                      Reaction
NAME        hexokinase
DEFINITION  ATP + D-Glucose <=> ADP + D-Glucose 6-phosphate
EQUATION    C00002 + C00031 <=> C00008 + C00092
ENZYME      2.7.1.1
PATHWAY     rn00010  Glycolysis / Gluconeogenesis
            rn00052  Galactose metabolism
///
\end{lstlisting}

\textbf{Key fields}: - \passthrough{\lstinline!DEFINITION!}:
Human-readable reaction equation - \passthrough{\lstinline!EQUATION!}:
KEGG compound IDs with stoichiometry - \passthrough{\lstinline!ENZYME!}:
EC number (Enzyme Commission classification)

\begin{center}\rule{0.5\linewidth}{0.5pt}\end{center}

\hypertarget{shypn-kegg-connector}{%
\section{9.3 SHYpn KEGG Connector}\label{shypn-kegg-connector}}

\hypertarget{keggconnector-class}{%
\subsection{9.3.1 KEGGConnector Class}\label{keggconnector-class}}

\textbf{Responsibilities}: - Send HTTP requests to KEGG REST API - Parse
plain-text responses - Cache results (avoid redundant API calls) -
Handle errors (network failures, invalid IDs)

\textbf{Implementation}:

\begin{lstlisting}[language=Python]
import requests
import re
from typing import Dict, Optional
from functools import lru_cache

class KEGGConnector:
    """Interface to KEGG REST API."""
    
    BASE_URL = "https://rest.kegg.jp"
    
    def __init__(self):
        self.session = requests.Session()
        self.session.headers.update({'User-Agent': 'SHYpn/1.0'})
    
    @lru_cache(maxsize=1000)
    def get_compound(self, compound_id: str) -> Optional[Dict[str, any]]:
        """Fetch compound details by KEGG ID.
        
        Args:
            compound_id: KEGG compound ID (e.g., "C00031")
        
        Returns:
            Dictionary with keys: name, formula, exact_mass, reactions, pathways
            Returns None if compound not found.
        """
        url = f"{self.BASE_URL}/get/{compound_id}"
        
        try:
            response = self.session.get(url, timeout=10)
            response.raise_for_status()
        except requests.RequestException as e:
            print(f"KEGG API error: {e}")
            return None
        
        # Parse plain-text response
        data = self._parse_compound_response(response.text)
        return data
    
    def _parse_compound_response(self, text: str) -> Dict[str, any]:
        """Parse KEGG compound response text."""
        data = {}
        
        # Extract name (first line after NAME)
        name_match = re.search(r'NAME\s+(.*)', text)
        if name_match:
            data['name'] = name_match.group(1).split(';')[0].strip()
        
        # Extract formula
        formula_match = re.search(r'FORMULA\s+(\S+)', text)
        if formula_match:
            data['formula'] = formula_match.group(1)
        
        # Extract exact mass
        mass_match = re.search(r'EXACT_MASS\s+([\d.]+)', text)
        if mass_match:
            data['exact_mass'] = float(mass_match.group(1))
        
        # Extract reactions (space-separated)
        reaction_match = re.search(r'REACTION\s+(.*)', text)
        if reaction_match:
            data['reactions'] = reaction_match.group(1).split()
        
        # Extract pathways
        pathway_matches = re.findall(r'PATHWAY\s+(map\d+)\s+(.*)', text)
        data['pathways'] = {pid: name.strip() for pid, name in pathway_matches}
        
        return data
    
    @lru_cache(maxsize=500)
    def get_reaction(self, reaction_id: str) -> Optional[Dict[str, any]]:
        """Fetch reaction details by KEGG ID.
        
        Args:
            reaction_id: KEGG reaction ID (e.g., "R00299")
        
        Returns:
            Dictionary with keys: name, definition, equation, enzyme, reversible
        """
        url = f"{self.BASE_URL}/get/{reaction_id}"
        
        try:
            response = self.session.get(url, timeout=10)
            response.raise_for_status()
        except requests.RequestException as e:
            print(f"KEGG API error: {e}")
            return None
        
        data = self._parse_reaction_response(response.text)
        return data
    
    def _parse_reaction_response(self, text: str) -> Dict[str, any]:
        """Parse KEGG reaction response text."""
        data = {}
        
        # Extract name
        name_match = re.search(r'NAME\s+(.*)', text)
        if name_match:
            data['name'] = name_match.group(1).strip()
        
        # Extract definition (human-readable equation)
        def_match = re.search(r'DEFINITION\s+(.*)', text)
        if def_match:
            data['definition'] = def_match.group(1).strip()
        
        # Extract equation (KEGG compound IDs)
        eq_match = re.search(r'EQUATION\s+(.*)', text)
        if eq_match:
            equation_str = eq_match.group(1).strip()
            data['equation'] = self._parse_equation(equation_str)
            data['reversible'] = '<=>' in equation_str
        
        # Extract EC number
        enzyme_match = re.search(r'ENZYME\s+([\d.]+)', text)
        if enzyme_match:
            data['ec_number'] = enzyme_match.group(1)
        
        return data
    
    def _parse_equation(self, equation_str: str) -> Dict[str, any]:
        """Parse KEGG equation string into structured data.
        
        Example:
            "C00002 + C00031 <=> C00008 + C00092"
            \textrightarrow
            {
                'substrates': [('C00002', 1), ('C00031', 1)],
                'products': [('C00008', 1), ('C00092', 1)]
            }
        """
        # Split by arrow (handle both => and <=>)
        if '<=>' in equation_str:
            left, right = equation_str.split('<=>')
        elif '=>' in equation_str:
            left, right = equation_str.split('=>')
        else:
            return {}
        
        def parse_side(side_str):
            """Parse one side of equation (e.g., "2 C00002 + C00031")."""
            compounds = []
            for term in side_str.strip().split('+'):
                term = term.strip()
                # Check for stoichiometric coefficient
                match = re.match(r'(\d+)\s+(\w+)', term)
                if match:
                    coeff = int(match.group(1))
                    compound_id = match.group(2)
                else:
                    coeff = 1
                    compound_id = term
                compounds.append((compound_id, coeff))
            return compounds
        
        return {
            'substrates': parse_side(left),
            'products': parse_side(right)
        }
    
    def search_compound(self, query: str) -> List[Dict[str, str]]:
        """Search compounds by name.
        
        Args:
            query: Search term (e.g., "glucose")
        
        Returns:
            List of matches: [{"id": "C00031", "name": "D-Glucose"}, ...]
        """
        url = f"{self.BASE_URL}/find/compound/{query}"
        
        try:
            response = self.session.get(url, timeout=10)
            response.raise_for_status()
        except requests.RequestException as e:
            print(f"KEGG API error: {e}")
            return []
        
        # Parse results (format: "cpd:C00031\tD-Glucose; Grape sugar")
        results = []
        for line in response.text.strip().split('\n'):
            if '\t' in line:
                compound_part, name_part = line.split('\t', 1)
                compound_id = compound_part.split(':')[1]
                name = name_part.split(';')[0].strip()
                results.append({'id': compound_id, 'name': name})
        
        return results
\end{lstlisting}

\hypertarget{usage-examples}{%
\subsection{9.3.2 Usage Examples}\label{usage-examples}}

\textbf{Example 1: Lookup compound by ID}:

\begin{lstlisting}[language=Python]
connector = KEGGConnector()

# Fetch glucose
glucose = connector.get_compound("C00031")
print(glucose)
# Output:
# {
#   'name': 'D-Glucose',
#   'formula': 'C6H12O6',
#   'exact_mass': 180.0634,
#   'reactions': ['R00299', 'R00300', ...],
#   'pathways': {'map00010': 'Glycolysis / Gluconeogenesis', ...}
# }
\end{lstlisting}

\textbf{Example 2: Search compound by name}:

\begin{lstlisting}[language=Python]
results = connector.search_compound("pyruvate")
print(results)
# Output:
# [
#   {'id': 'C00022', 'name': 'Pyruvate'},
#   {'id': 'C00074', 'name': 'Phosphoenolpyruvate'},
#   ...
# ]
\end{lstlisting}

\textbf{Example 3: Fetch reaction}:

\begin{lstlisting}[language=Python]
reaction = connector.get_reaction("R00299")
print(reaction)
# Output:
# {
#   'name': 'hexokinase',
#   'definition': 'ATP + D-Glucose <=> ADP + D-Glucose 6-phosphate',
#   'equation': {
#     'substrates': [('C00002', 1), ('C00031', 1)],
#     'products': [('C00008', 1), ('C00092', 1)]
#   },
#   'ec_number': '2.7.1.1',
#   'reversible': True
# }
\end{lstlisting}

\begin{center}\rule{0.5\linewidth}{0.5pt}\end{center}

\hypertarget{automatic-model-enrichment}{%
\section{9.4 Automatic Model
Enrichment}\label{automatic-model-enrichment}}

\hypertarget{place-formula-auto-fill}{%
\subsection{9.4.1 Place Formula
Auto-Fill}\label{place-formula-auto-fill}}

\textbf{User workflow}: 1. User creates a place and enters KEGG ID
(e.g., ``C00031'') 2. User clicks \textbf{``Fetch from KEGG''} button 3.
SHYpn fetches formula, name, and fills in form fields

\textbf{Implementation}:

\begin{lstlisting}[language=Python]
class PlacePropertyPanel:
    """GTK panel for editing place properties."""
    
    def on_fetch_from_kegg_clicked(self, button):
        """Handle "Fetch from KEGG" button click."""
        kegg_id = self.kegg_id_entry.get_text().strip()
        
        if not kegg_id:
            self.show_error("Please enter a KEGG ID (e.g., C00031)")
            return
        
        # Show loading spinner
        self.set_loading(True)
        
        # Fetch from KEGG (asynchronous to avoid blocking UI)
        connector = KEGGConnector()
        compound = connector.get_compound(kegg_id)
        
        self.set_loading(False)
        
        if compound is None:
            self.show_error(f"KEGG compound {kegg_id} not found")
            return
        
        # Auto-fill form fields
        if 'name' in compound:
            self.name_entry.set_text(compound['name'])
        
        if 'formula' in compound:
            self.formula_entry.set_text(compound['formula'])
            # Parse and validate formula
            try:
                parsed_formula = BiochemicalFormula.parse(compound['formula'])
                self.formula_valid_icon.set_from_icon_name('emblem-ok-symbolic')
            except ValueError as e:
                self.show_warning(f"Formula parsing error: {e}")
        
        self.show_info(f"Fetched {compound['name']} from KEGG")
\end{lstlisting}

\textbf{Benefits}: - \textbf{Speed}: No manual formula entry -
\textbf{Accuracy}: KEGG formulas are curated - \textbf{Consistency}:
Same compound always has same formula

\hypertarget{reaction-import}{%
\subsection{9.4.2 Reaction Import}\label{reaction-import}}

\textbf{User workflow}: 1. User selects \textbf{``Import from KEGG''}
menu item 2. Dialog prompts for KEGG reaction ID (e.g., ``R00299'') 3.
SHYpn fetches reaction, creates: - \textbf{Places} for all substrates
and products (if not already present) - \textbf{Transition} with
appropriate name and EC number - \textbf{Arcs} with stoichiometric
weights 4. User reviews and adjusts parameters (Vmax, Km)

\textbf{Implementation}:

\begin{lstlisting}[language=Python]
class ModelManager:
    """Manages Bio-PN model state."""
    
    def import_kegg_reaction(self, reaction_id: str) -> None:
        """Import a KEGG reaction into the model.
        
        Creates:
        - Places for substrates and products (auto-fetch formulas)
        - Transition (name from KEGG)
        - Normal arcs with stoichiometric weights
        """
        connector = KEGGConnector()
        
        # Fetch reaction
        reaction_data = connector.get_reaction(reaction_id)
        if reaction_data is None:
            raise ValueError(f"KEGG reaction {reaction_id} not found")
        
        # Create transition
        transition = Transition(
            id=f"t_{reaction_id}",
            name=reaction_data.get('name', reaction_id),
            transition_type=TransitionType.CONTINUOUS,
            rate_function=MassActionRate(k=1.0, substrate_ids=[]),  # Placeholder
            reversible=reaction_data.get('reversible', False),
            ec_number=reaction_data.get('ec_number'),
            kegg_reaction_id=reaction_id
        )
        self.add_transition(transition)
        
        # Create places for substrates
        equation = reaction_data.get('equation', {})
        for compound_id, coeff in equation.get('substrates', []):
            place_id = f"p_{compound_id}"
            
            # Check if place already exists
            if place_id not in self.model.places:
                # Fetch compound data
                compound_data = connector.get_compound(compound_id)
                place = Place(
                    id=place_id,
                    name=compound_data.get('name', compound_id),
                    formula=BiochemicalFormula.parse(compound_data.get('formula', '')),
                    initial_marking=1.0,  # Default
                    kegg_id=compound_id
                )
                self.add_place(place)
            
            # Create input arc (place \textrightarrow transition)
            arc = Arc(
                id=f"a_{place_id}_{transition.id}",
                source=place_id,
                target=transition.id,
                arc_type=ArcType.NORMAL,
                weight=coeff
            )
            self.add_arc(arc)
        
        # Create places for products
        for compound_id, coeff in equation.get('products', []):
            place_id = f"p_{compound_id}"
            
            if place_id not in self.model.places:
                compound_data = connector.get_compound(compound_id)
                place = Place(
                    id=place_id,
                    name=compound_data.get('name', compound_id),
                    formula=BiochemicalFormula.parse(compound_data.get('formula', '')),
                    initial_marking=0.0,  # Product starts at 0
                    kegg_id=compound_id
                )
                self.add_place(place)
            
            # Create output arc (transition \textrightarrow place)
            arc = Arc(
                id=f"a_{transition.id}_{place_id}",
                source=transition.id,
                target=place_id,
                arc_type=ArcType.NORMAL,
                weight=coeff
            )
            self.add_arc(arc)
        
        # Verify elemental balance
        balance = self.validate_elemental_balance(transition)
        if balance:
            print(f"Warning: Reaction {reaction_id} is not elementally balanced:")
            print(f"  Imbalance: {balance}")
\end{lstlisting}

\textbf{Example}: Importing hexokinase (R00299):

\begin{lstlisting}
Before import: Empty model

After import:
  Places:
    - p_C00002 (ATP, C10H16N5O13P3)
    - p_C00031 (D-Glucose, C6H12O6)
    - p_C00008 (ADP, C10H15N5O10P2)
    - p_C00092 (D-Glucose 6-phosphate, C6H13O9P)
  
  Transition:
    - t_R00299 (hexokinase, EC 2.7.1.1, Continuous)
  
  Arcs:
    - p_C00002 \textrightarrow t_R00299 (weight=1)
    - p_C00031 \textrightarrow t_R00299 (weight=1)
    - t_R00299 \textrightarrow p_C00008 (weight=1)
    - t_R00299 \textrightarrow p_C00092 (weight=1)
\end{lstlisting}

\hypertarget{pathway-import}{%
\subsection{9.4.3 Pathway Import}\label{pathway-import}}

\textbf{User workflow}: 1. User selects \textbf{``Import KEGG Pathway''}
menu item 2. Dialog shows list of pathways (e.g., ``Glycolysis /
Gluconeogenesis (map00010)'') 3. User selects pathway 4. SHYpn fetches
all reactions in pathway, imports each 5. Automatic layout applied
(hierarchical, left-to-right)

\textbf{Implementation} (simplified):

\begin{lstlisting}[language=Python]
def import_kegg_pathway(self, pathway_id: str) -> None:
    """Import entire KEGG pathway.
    
    Args:
        pathway_id: KEGG pathway ID (e.g., "map00010")
    """
    connector = KEGGConnector()
    
    # Fetch pathway (returns list of reaction IDs)
    url = f"{connector.BASE_URL}/link/reaction/{pathway_id}"
    response = requests.get(url)
    reaction_ids = [line.split('\t')[1].split(':')[1] 
                    for line in response.text.strip().split('\n')]
    
    # Import each reaction
    for reaction_id in reaction_ids:
        try:
            self.import_kegg_reaction(reaction_id)
        except Exception as e:
            print(f"Failed to import {reaction_id}: {e}")
    
    # Auto-layout
    self.apply_hierarchical_layout()
\end{lstlisting}

\textbf{Example}: Importing glycolysis (map00010): - \textbf{10
reactions} imported (R00299, R00300, \ldots, R00658) - \textbf{13
places} created (glucose, G6P, F6P, \ldots, pyruvate) -
\textbf{Execution time}: \textasciitilde4 seconds (10 API calls)

\begin{center}\rule{0.5\linewidth}{0.5pt}\end{center}

\hypertarget{cofactor-suggestion}{%
\section{9.5 Cofactor Suggestion}\label{cofactor-suggestion}}

\hypertarget{problem-missing-cofactors}{%
\subsection{9.5.1 Problem: Missing
Cofactors}\label{problem-missing-cofactors}}

\textbf{KEGG reaction equations often omit common cofactors}: - Water
(H₂O) added in hydrolysis - Protons (H⁺) released in oxidation -
Inorganic phosphate (Pi) in ATP hydrolysis

\textbf{Example}: KEGG R00299 (hexokinase) - \textbf{KEGG equation}:
\passthrough{\lstinline!ATP + Glucose <=> ADP + G6P!} - \textbf{Actual
reaction}: \passthrough{\lstinline!ATP + Glucose <=> ADP + G6P + H⁺!} -
\textbf{Elemental imbalance} without H⁺: Left has 16 H, right has 15 H

\textbf{SHYpn approach}: \textbf{Automatically suggest missing
cofactors} based on elemental imbalance.

\hypertarget{cofactor-suggestion-algorithm}{%
\subsection{9.5.2 Cofactor Suggestion
Algorithm}\label{cofactor-suggestion-algorithm}}

\textbf{Algorithm}:

\begin{lstlisting}
Input: Transition t with unbalanced reaction
Output: List of suggested cofactors with stoichiometry

1. Compute elemental imbalance Δ = Input_elements - Output_elements
   (e.g., Δ = {"H": 1, "O": 0, "C": 0, ...})

2. Define common cofactor formulas:
   - H₂O: {"H": 2, "O": 1}
   - H⁺: {"H": 1}
   - HPO₄²⁻ (phosphate): {"H": 1, "P": 1, "O": 4}
   - NAD⁺: {"C": 21, "H": 27, "N": 7, "O": 14, "P": 2}
   - NADH: {"C": 21, "H": 29, "N": 7, "O": 14, "P": 2}
   - CO₂: {"C": 1, "O": 2}

3. Try combinations of cofactors (linear Diophantine equation):
   Find coefficients (n₁, n₂, ...) such that:
     Δ = n₁ · Formula₁ + n₂ · Formula₂ + ...
   
4. Rank suggestions by:
   - Fewest cofactors (prefer simple)
   - Biochemical plausibility (e.g., H⁺ + H₂O more likely than NAD⁺ alone)

5. Return top 3 suggestions

Example:
  Imbalance: {"H": 1}
  Suggestion 1: +1 H⁺ (rank=1.0)
  Suggestion 2: +0.5 H₂O, -0.5 O (rank=0.3)  [less plausible]
\end{lstlisting}

\textbf{Implementation}:

\begin{lstlisting}[language=Python]
class CofactorSuggester:
    """Suggests missing cofactors for unbalanced reactions."""
    
    COMMON_COFACTORS = {
        'H+': BiochemicalFormula.parse('H'),
        'H2O': BiochemicalFormula.parse('H2O'),
        'Pi': BiochemicalFormula.parse('HPO4'),
        'NAD+': BiochemicalFormula.parse('C21H27N7O14P2'),
        'NADH': BiochemicalFormula.parse('C21H29N7O14P2'),
        'CO2': BiochemicalFormula.parse('CO2'),
    }
    
    def suggest_cofactors(self, transition: Transition, model: BioPetriNet) -> List[Dict]:
        """Suggest cofactors to balance reaction.
        
        Returns:
            List of suggestions: [
              {'cofactors': {'H+': 1}, 'score': 1.0},
              ...
            ]
        """
        # Compute elemental imbalance
        imbalance = self.compute_imbalance(transition, model)
        
        if not imbalance:
            return []  # Already balanced
        
        suggestions = []
        
        # Try single cofactors
        for name, formula in self.COMMON_COFACTORS.items():
            if self.can_balance(imbalance, formula):
                coeff = self.compute_coefficient(imbalance, formula)
                suggestions.append({
                    'cofactors': {name: coeff},
                    'score': 1.0
                })
        
        # Try pairs of cofactors (for complex imbalances)
        for name1, formula1 in self.COMMON_COFACTORS.items():
            for name2, formula2 in self.COMMON_COFACTORS.items():
                if name1 >= name2:
                    continue  # Avoid duplicates
                
                coeffs = self.solve_pair(imbalance, formula1, formula2)
                if coeffs:
                    suggestions.append({
                        'cofactors': {name1: coeffs[0], name2: coeffs[1]},
                        'score': 0.5  # Prefer single cofactor
                    })
        
        # Sort by score
        suggestions.sort(key=lambda s: s['score'], reverse=True)
        return suggestions[:3]  # Top 3
    
    def compute_imbalance(self, transition: Transition, model: BioPetriNet) -> Dict[str, int]:
        """Compute elemental imbalance (Inputs - Outputs)."""
        input_elements = {}
        output_elements = {}
        
        # Sum input place formulas (weighted by arc weight)
        for arc in model.arcs:
            if arc.target == transition.id and arc.arc_type == ArcType.NORMAL:
                place = model.places[arc.source]
                for element, count in place.formula.elements.items():
                    input_elements[element] = input_elements.get(element, 0) + count * arc.weight
        
        # Sum output place formulas
        for arc in model.arcs:
            if arc.source == transition.id and arc.arc_type == ArcType.NORMAL:
                place = model.places[arc.target]
                for element, count in place.formula.elements.items():
                    output_elements[element] = output_elements.get(element, 0) + count * arc.weight
        
        # Compute difference
        imbalance = {}
        all_elements = set(input_elements.keys()) | set(output_elements.keys())
        for element in all_elements:
            diff = input_elements.get(element, 0) - output_elements.get(element, 0)
            if diff != 0:
                imbalance[element] = diff
        
        return imbalance
\end{lstlisting}

\hypertarget{example-hexokinase-balance}{%
\subsection{9.5.3 Example: Hexokinase
Balance}\label{example-hexokinase-balance}}

\textbf{Initial reaction} (from KEGG R00299):

\begin{lstlisting}
ATP (C10H16N5O13P3) + Glucose (C6H12O6) 
  \textrightarrow ADP (C10H15N5O10P2) + G6P (C6H13O9P)
\end{lstlisting}

\textbf{Elemental count}: - \textbf{Inputs}: C=16, H=28, N=5, O=19, P=4
- \textbf{Outputs}: C=16, H=28, N=5, O=19, P=3 - \textbf{Imbalance}:
P=+1

\textbf{Cofactor suggester}:

\begin{lstlisting}[language=Python]
suggester = CofactorSuggester()
suggestions = suggester.suggest_cofactors(t_hexokinase, model)

print(suggestions)
# Output:
# [
#   {'cofactors': {'Pi': 1}, 'score': 1.0},
#   {'cofactors': {'H+': 1, 'Pi': 1}, 'score': 0.5}  # Over-correction
# ]
\end{lstlisting}

\textbf{User action}: Accept suggestion \textrightarrow Add Pi place and output arc

\textbf{Corrected reaction}:

\begin{lstlisting}
ATP + Glucose \textrightarrow ADP + G6P + Pi
\end{lstlisting}

\begin{center}\rule{0.5\linewidth}{0.5pt}\end{center}

\hypertarget{integration-with-model-validation}{%
\section{9.6 Integration with Model
Validation}\label{integration-with-model-validation}}

\hypertarget{validation-workflow}{%
\subsection{9.6.1 Validation Workflow}\label{validation-workflow}}

\textbf{SHYpn validation steps} (triggered on model save):

\begin{enumerate}
\def\labelenumi{\arabic{enumi}.}
\tightlist
\item
  \textbf{Structural validation} (well-formedness constraints C1-C8)
\item
  \textbf{Elemental balance check} (for all transitions)
\item
  \textbf{If imbalance detected}:

  \begin{itemize}
  \tightlist
  \item
    Show warning dialog
  \item
    Display suggested cofactors
  \item
    User can accept, modify, or ignore
  \end{itemize}
\end{enumerate}

\textbf{Implementation}:

\begin{lstlisting}[language=Python]
class ModelValidator:
    """Validates Bio-PN model correctness."""
    
    def validate_model(self, model: BioPetriNet) -> List[ValidationError]:
        """Run all validation checks.
        
        Returns:
            List of errors/warnings
        """
        errors = []
        
        # Check well-formedness
        errors.extend(self.check_well_formedness(model))
        
        # Check elemental balance for all transitions
        for transition in model.transitions.values():
            imbalance = self.check_elemental_balance(transition, model)
            if imbalance:
                # Generate cofactor suggestions
                suggester = CofactorSuggester()
                suggestions = suggester.suggest_cofactors(transition, model)
                
                errors.append(ValidationError(
                    level='WARNING',
                    message=f"Transition {transition.name} is not elementally balanced",
                    details=f"Imbalance: {imbalance}",
                    suggestions=suggestions
                ))
        
        return errors
\end{lstlisting}

\textbf{User experience}: - \textbf{Green checkmark}: Model fully
validated (structurally sound + elementally balanced) - \textbf{Yellow
warning}: Model has imbalances, but suggestions available - \textbf{Red
error}: Structural issues (e.g., test arc with weight ≠ 1)

\begin{center}\rule{0.5\linewidth}{0.5pt}\end{center}

\hypertarget{performance-and-caching}{%
\section{9.7 Performance and Caching}\label{performance-and-caching}}

\hypertarget{api-call-minimization}{%
\subsection{9.7.1 API Call Minimization}\label{api-call-minimization}}

\textbf{Challenge}: KEGG REST API is relatively slow (200-500 ms per
request)

\textbf{Strategies}: 1. \textbf{LRU cache}:
\passthrough{\lstinline!@lru\_cache!} decorator on
\passthrough{\lstinline!get\_compound()!} and
\passthrough{\lstinline!get\_reaction()!} - Avoids repeated fetches of
same compound - 1000-entry cache (sufficient for typical models) 2.
\textbf{Batch requests}: When importing pathway, fetch all reactions in
parallel 3. \textbf{Offline mode}: Store downloaded compounds in local
SQLite database - First fetch: Download from KEGG + save to DB -
Subsequent: Read from DB (10× faster)

\textbf{Benchmark} (import glycolysis pathway): - \textbf{Without
cache}: 10 reactions × 4 compounds × 500 ms = \textbf{20 seconds} -
\textbf{With cache}: 10 reactions × 1 fetch + 13 compounds × 500 ms =
\textbf{7 seconds} - \textbf{With offline DB}: 10 reactions + 13
compounds from DB = \textbf{0.3 seconds}

\hypertarget{offline-database}{%
\subsection{9.7.2 Offline Database}\label{offline-database}}

\textbf{Schema} (SQLite):

\begin{lstlisting}[language=SQL]
CREATE TABLE compounds (
    kegg_id TEXT PRIMARY KEY,
    name TEXT,
    formula TEXT,
    exact_mass REAL,
    fetched_at TIMESTAMP DEFAULT CURRENT_TIMESTAMP
);

CREATE TABLE reactions (
    kegg_id TEXT PRIMARY KEY,
    name TEXT,
    definition TEXT,
    equation TEXT,  -- JSON string
    ec_number TEXT,
    reversible BOOLEAN,
    fetched_at TIMESTAMP DEFAULT CURRENT_TIMESTAMP
);
\end{lstlisting}

\textbf{Implementation}:

\begin{lstlisting}[language=Python]
import sqlite3
import json
from datetime import datetime, timedelta

class KEGGDatabase:
    """Local cache of KEGG data."""
    
    def __init__(self, db_path='kegg_cache.db'):
        self.conn = sqlite3.connect(db_path)
        self.conn.row_factory = sqlite3.Row
        self.create_tables()
    
    def get_compound(self, compound_id: str, max_age_days=30) -> Optional[Dict]:
        """Fetch compound from local DB (with freshness check)."""
        cursor = self.conn.execute(
            'SELECT * FROM compounds WHERE kegg_id = ?', (compound_id,)
        )
        row = cursor.fetchone()
        
        if row is None:
            return None
        
        # Check freshness
        fetched_at = datetime.fromisoformat(row['fetched_at'])
        if datetime.now() - fetched_at > timedelta(days=max_age_days):
            return None  # Stale, refetch from KEGG
        
        return dict(row)
    
    def save_compound(self, compound_id: str, data: Dict) -> None:
        """Save compound to local DB."""
        self.conn.execute('''
            INSERT OR REPLACE INTO compounds (kegg_id, name, formula, exact_mass)
            VALUES (?, ?, ?, ?)
        ''', (compound_id, data.get('name'), data.get('formula'), data.get('exact_mass')))
        self.conn.commit()
\end{lstlisting}

\textbf{Hybrid strategy}: Check local DB first, fallback to KEGG API

\begin{center}\rule{0.5\linewidth}{0.5pt}\end{center}

\hypertarget{limitations-and-future-work}{%
\section{9.8 Limitations and Future
Work}\label{limitations-and-future-work}}

\hypertarget{current-limitations}{%
\subsection{9.8.1 Current Limitations}\label{current-limitations}}

\begin{enumerate}
\def\labelenumi{\arabic{enumi}.}
\tightlist
\item
  \textbf{Protonation states}: KEGG formulas are at pH 7, but
  protonation varies by pH

  \begin{itemize}
  \tightlist
  \item
    \textbf{Impact}: Elemental counts slightly off (H⁺ count ambiguous)
  \item
    \textbf{Workaround}: Assume pH 7, ignore small H imbalances
  \end{itemize}
\item
  \textbf{Incomplete stoichiometry}: Some KEGG reactions omit minor
  products

  \begin{itemize}
  \tightlist
  \item
    \textbf{Example}: R00299 omits H⁺ released by ATP hydrolysis
  \item
    \textbf{Solution}: Cofactor suggester partially mitigates
  \end{itemize}
\item
  \textbf{Rate constant data}: KEGG does not provide kinetic parameters
  (Vmax, Km)

  \begin{itemize}
  \tightlist
  \item
    \textbf{Solution}: Chapter 10 (BRENDA integration for parameters)
  \end{itemize}
\item
  \textbf{Pathway topology}: KEGG pathways are lists of reactions (no
  explicit network structure)

  \begin{itemize}
  \tightlist
  \item
    \textbf{Solution}: SHYpn infers connections from shared metabolites
  \end{itemize}
\end{enumerate}

\hypertarget{future-enhancements}{%
\subsection{9.8.2 Future Enhancements}\label{future-enhancements}}

\begin{enumerate}
\def\labelenumi{\arabic{enumi}.}
\tightlist
\item
  \textbf{ChEBI integration}: Alternative database with more detailed
  formulas
\item
  \textbf{MetaCyc integration}: Curated metabolic pathways with better
  stoichiometry
\item
  \textbf{Rhea integration}: Reaction database with balanced equations
\item
  \textbf{Automatic rate constant estimation}: Use machine learning to
  predict Km from enzyme structure
\end{enumerate}

\begin{center}\rule{0.5\linewidth}{0.5pt}\end{center}

\hypertarget{summary}{%
\section{9.9 Summary}\label{summary}}

\textbf{Chapter 9 presented KEGG integration for automatic model
enrichment}:

\begin{enumerate}
\def\labelenumi{\arabic{enumi}.}
\tightlist
\item
  \textbf{KEGG REST API}: Fetches compound formulas and reaction
  stoichiometry
\item
  \textbf{KEGGConnector class}: Python wrapper with caching, error
  handling
\item
  \textbf{Automatic enrichment}:

  \begin{itemize}
  \tightlist
  \item
    Place formula auto-fill (click ``Fetch from KEGG'')
  \item
    Reaction import (creates places, transition, arcs)
  \item
    Pathway import (entire pathways in seconds)
  \end{itemize}
\item
  \textbf{Cofactor suggestion}: Algorithm suggests missing H₂O, H⁺, Pi
  based on elemental imbalance
\item
  \textbf{Performance}: LRU cache + offline SQLite DB \textrightarrow 60× speedup
\item
  \textbf{Validation integration}: Elemental balance checks with
  actionable suggestions
\end{enumerate}

\textbf{Key benefits}: - \textbf{Speed}: Import glycolysis (10
reactions) in \textless5 seconds - \textbf{Accuracy}: KEGG formulas are
manually curated (95\%+ correct) - \textbf{Completeness}: Cofactor
suggester fills gaps in KEGG data

\textbf{Example}: Complete upper glycolysis pathway (glucose \textrightarrow pyruvate)
created in 30 seconds with 0 manual formula entries.

\textbf{Next chapter} (Chapter 10): BRENDA integration for kinetic
parameter inference (Vmax, Km).
